在均匀、无耗散的介质中，波数为 \(k_0>0\) 的二维标量单色光场写为 \(E(x,z)=\psi(x,z)e^{ik_0z}\)。假设傍轴近似适用，选横向标尺 \(x_0>0\)，定义 \(s=x/x_0\)、\(\xi=z/(k_0x_0^2)\)，因而包络满足
\[
i\psi_\xi+\frac12\psi_{ss}=0.
\]
考虑 \(a>0\) 的截断 Airy 初始包络 \(\psi(s,0)=\operatorname{Ai}(s)e^{as}\)。取空间傅里叶约定 \(\widehat\psi(k,\xi)=\int_{\mathbb R}\psi(s,\xi)e^{-iks}\,ds\)、逆变换为 \((2\pi)^{-1}\int_{\mathbb R}\widehat\psi e^{iks}dk\)。所有功率均按单位垂直纸面的长度并舍去共同材料常数。
(1) 从亥姆霍兹方程导出该无量纲傍轴方程；推导局部功率连续性方程及守恒总功率，写出横向光流 \(j(s,\xi)\) 的表达式与边界积分条件。
(2) 利用 Airy 函数的积分表示和傅里叶传播因子，严格计算任意 \(\xi\) 的复包络 \(\psi(s,\xi)\)。代入偏微分方程复核，说明 \(a\to0^+\) 时可得到哪一个理想但不可归一化的弯曲光束。
(3) 计算 \(|\widehat\psi(k,0)|^2\)、总功率 \(\mathcal P\)、平均横向波数 \(\langle k\rangle\)、质心 \(\langle s\rangle(\xi)\) 和方差 \(\operatorname{Var}(s)(\xi)\)。证明虽然理想 Airy 主亮瓣沿抛物线移动，任何固定 \(a>0\) 的有限功率光束质心却作匀速运动（本题为静止）。
(4) 计算任意 \(a>0\) 时的局部光流 \(j\)（允许用 Airy 函数对复变量的导数表示），证明其横向积分为零；在 \(a=0\) 的理想极限写出主亮瓣轨迹的物理尺寸形式与形式上的局部流速，并用谱相位的驻点合并条件说明截断后的“加速”是亮纹结构而非整束光的横向外力加速。